\documentclass[11pt,a4paper]{article}

\usepackage[margin=25mm]{geometry}
\usepackage{fontspec}
\usepackage{libertinus}
\newfontfamily\devanagarifont{Kohinoor Devanagari}
\newfontfamily\banglafont{Kohinoor Bangla}
\usepackage{microtype}
\usepackage{booktabs}
\usepackage{tabularx}
\usepackage{siunitx}
\usepackage{amsmath}
\usepackage{graphicx}
\usepackage{pgfplots}
\usepackage[numbers,sort&compress]{natbib}
\usepackage[hidelinks]{hyperref}
\usepackage[nameinlink,noabbrev]{cleveref}
\usepackage{xcolor}
\usepackage{enumitem}
\usepackage{caption}
\usepackage{float}

\pgfplotsset{compat=1.18}
\sisetup{detect-all}
\definecolor{accent}{HTML}{315E8A}
\definecolor{accentlight}{HTML}{8BB8D8}
\captionsetup{font=small,labelfont=bf}
\setlist{nosep,leftmargin=*}
\setlength{\parindent}{1.2em}
\setlength{\parskip}{0pt}

\newcommand{\QualificationCorrect}{5/5}
\newcommand{\PilotCorrect}{47/60}
\newcommand{\PilotAccuracy}{78.3\%}
\newcommand{\PilotMeanLatency}{252.5}
\newcommand{\PilotMedianLatency}{194.9}
\newcommand{\PilotMeanTokens}{46944}
\newcommand{\PilotTruncations}{2}
\newcommand{\PilotOperationalFailures}{0}
\newcommand{\MainPartOneCorrect}{69/101}
\newcommand{\MainPartOneAccuracy}{68.3\%}
\newcommand{\MainPartOneMeanLatency}{278.2}
\newcommand{\MainPartOneMedianLatency}{222.5}
\newcommand{\MainPartOneMeanTokens}{51611}
\newcommand{\MainPartOneTruncations}{6}
\newcommand{\MainPartOneOperationalFailures}{0}


\title{Beyond Arabic Digits: Evaluating Sudoku Solving Under Symbol-Set Variation}
\author{Arnav Bharti}
\date{September 2026}

\begin{document}
\maketitle

\begin{abstract}
Sudoku has the same solution when its nine values are consistently renamed. This makes it a useful test of whether a language model follows the puzzle's constraints across different written forms. We test this question by holding each Sudoku puzzle fixed while changing only its value symbols. We generated 300 unique-solution puzzles and evaluated GPT-OSS-120B on a fixed sample of 60 puzzles, each written with nine different symbol sets. A completed answer was correct only when it preserved every clue and satisfied every Sudoku constraint. The model solved \MainCorrect{} requests (\MainAccuracy{}) without operational failures. Accuracy was 75.0\% with Arabic digits and uppercase Latin letters, compared with 61.7\% with Greek letters. Follow-up tests found perfect performance on 75 elementary symbol tasks, but no single tested factor explained the Sudoku differences. The results show that a model's success on a logical puzzle can depend on how its values are written.
\end{abstract}

\section{Introduction}
Large language models can produce fluent answers to problems that require several steps of reasoning. Their accuracy can also change when a problem is reworded without changing its answer. This raises a basic question for reasoning benchmarks. Does a model solve the underlying problem, or does its success depend on familiar ways of writing that problem?

Sudoku makes the question testable. Every valid solution must preserve the given cells and contain each value exactly once in every row, column, and $3\times3$ box. We can replace the values 1 through 9 with nine other distinct symbols while leaving every logical constraint unchanged. The resulting puzzles have the same clues, solution, and difficulty. An accuracy difference between these versions therefore measures sensitivity to the representation, although it does not by itself identify the cause.

Prior work has evaluated language models on Sudoku and other reasoning tasks \citep{seely2025sudokubench,shah2024causal}. Other studies have shown that controlled changes to mathematical questions or symbol assignments can affect model behaviour \citep{mirzadeh2025gsmsymbolic,wei2023symbol}. These findings motivate a paired test in which every representation uses the same Sudoku puzzles and an exact checker scores only the final grid.

We ask whether GPT-OSS-120B retains Sudoku accuracy when the values are renamed. We first construct and audit a dataset of 300 puzzles. We then evaluate 60 held-out puzzles in nine representations, giving 540 model requests. Follow-up experiments test input and output mapping, label length, familiar but conflicting meanings, prompt wording, output format, and answer revision. The model solved 75.0\% of requests with Arabic digits and 61.7\% with Greek letters. The difference motivates the mechanism tests, but their small samples do not support a single causal explanation.

This paper contributes an exactly scored, paired representation benchmark, a frozen model and request protocol, and controlled follow-up experiments that separate several possible sources of representation sensitivity. The study measures complete-grid accuracy. It does not score the model's internal reasoning text.

\section{Related work and background}
\subsection{Sudoku as an exactly verifiable constraint problem}
Let $G=\{g_{r,c}\mid r,c\in\{1,\ldots,9\}\}$ denote the 81 grid cells and let $D=\{1,\ldots,9\}$ be the value domain. A completed grid is correct only if it preserves every clue and, for every row $r$, column $c$, and box $b$,
\begin{align}
\{g_{r,1},\ldots,g_{r,9}\}&=D,\\
\{g_{1,c},\ldots,g_{9,c}\}&=D,\\
\{g_{r,c}:(r,c)\in b\}&=D.
\end{align}
This definition permits deterministic evaluation of all cells and constraints. Sudoku difficulty depends on the deductions needed to solve a puzzle, not just its clue count \citep{pelanek2014sudoku}. Existing Sudoku evaluations test model performance on puzzles and variants \citep{seely2025sudokubench}. Solver-trace training has also been used to elicit search on logic puzzles \citep{shah2024causal}.

Our study uses the same Sudoku instance under several value alphabets. This paired design isolates the change in visible labels from changes in puzzle content. We score only the completed grid, so the evaluation does not require judging the quality of a reasoning trace.

\subsection{Representation invariance}
Define abstract values $V=\{V_1,\ldots,V_9\}$ and a visible alphabet $S=\{s_1,\ldots,s_9\}$. A representation is a bijection $\sigma:V\leftrightarrow S$. If $p$ is a Sudoku puzzle and $f$ is the model-and-parser pipeline, exact invariance requires
\begin{equation}
\sigma^{-1}\!\left(f(\sigma(p))\right)=f(p)
\end{equation}
whenever both sides return completed grids. The output must also pass the Sudoku checker. A model could consistently produce the same invalid grid under several encodings. Controlled changes to mathematical questions have exposed related sensitivities in language-model reasoning \citep{mirzadeh2025gsmsymbolic}. Our transformation keeps the full Sudoku puzzle fixed and changes its value labels.

\subsection{Binding, tokenization, and verification}
Symbol tuning can improve a model's use of arbitrary labels \citep{wei2023symbol}. Yet an 81-cell Sudoku answer requires the chosen mapping to remain consistent across the whole grid. Digit-renaming tests in arithmetic have reported related failures \citep{yan2025addition}. Those tests do not show which stage fails in a Sudoku answer.

Tokenization divides text into model input units called tokens. Different labels can require different numbers of tokens, and tokenizer choices can affect language-model performance \citep{bostrom2020bpe,rust2021tokenizer}. We therefore measure label cost and test neutral labels with controlled token lengths. We also use an exact checker to assess answer revision, since a model's own statement that an answer is correct is not a validity test \citep{stechly2023wrong}.

\section{Methodology}
\subsection{Study logic}
The study asks whether Sudoku solving is invariant to a change in visible value labels.  The abstract values are $V=\{V_1,\ldots,V_9\}$.  A representation is a bijection from these values to nine visible symbols.  Applying a bijection changes how the puzzle looks, but it does not change any clue position, row, column, box, or solution.  Therefore, a difference in accuracy between two representations cannot be caused by a different logical puzzle.

The experiments are ordered around a possible failure chain:
\begin{equation}
\text{decode the input}\;\longrightarrow\;\text{bind labels to values and solve}\;\longrightarrow\;\text{encode the output}.
\end{equation}
The main benchmark first measures the total representation effect.  Later experiments then change one part of the pipeline at a time.  They test input and output separately, token length, arbitrary mappings, semantic conflict, prompt and format choices, and answer revision.  This order is important: the main benchmark can show that a representation effect exists, while the mechanism experiments are needed to explain where it comes from.

Three research questions organize the results. RQ1 asks how complete-grid accuracy changes across value alphabets and difficulty tiers. RQ2 asks which tested features of a representation account for those changes. RQ3 asks whether prompt changes or answer revision improve correctness on the selected diagnostic puzzles. These questions concern observed model behaviour under the specified prompts and samples. They do not assume that any one feature has an isolated causal effect.

\subsection{Shared experimental controls}
\subsubsection{Model and inference}
The reported experiments use the open-weight GPT-OSS-120B checkpoint \citep{openai2025gptoss}.  Its repository identifier is \texttt{openai/\allowbreak gpt-oss-120b}, and the pinned revision is \texttt{b5c939de\allowbreak8f754692\allowbreak c1647ca7\allowbreak9fbf85e8\allowbreak c1e70f8a}.  Inference runs through vLLM on one NVIDIA H100 GPU with 12 CPU cores and 96~GB of host memory.  The chat template requests high reasoning effort.  Sampling uses temperature 1 and top-$p$ 1, with master seed 20260826.  The model context and maximum generation length are 131,072 tokens.

GPT-OSS returns Harmony analysis and final channels.  The backend stores them separately, and only the final channel is evaluated.  The visible reasoning trace is not an outcome of this thesis.  Scheduler waiting time is also excluded from latency. Latency measures model generation for one request.

\subsubsection{Open-weight model screening and selection}
The final benchmark model was selected through staged local screening rather than assumed in advance.  Table~\ref{tab:model-screening} separates completed inference evidence from models that were only prepared for execution.  This distinction is necessary because installation, configuration, or a pending scheduler job does not constitute a model result.

\begin{table}[H]
\centering
\footnotesize
\caption{Open-weight model selection. Prepared candidates have no scored Sudoku screen.}
\label{tab:model-screening}
\begin{tabularx}{\linewidth}{@{}p{0.34\linewidth}X@{}}
\toprule
Model & Completed evidence and decision \\
\midrule
NVIDIA Llama-3.3 Nemotron Super 49B v1.5 FP8 & Three answer-only settings scored 0/5, 0/15, and 0/15. Constrained output fixed formatting but not correctness, so we did not advance it. \\
Qwen3.8-27B (earlier non-FP8 screen) & One easy and one hard puzzle were correct. The medium request had no final grid. We did not advance this three-item screen. \\
GPT-OSS-120B & All three timing puzzles and all five qualification puzzles were correct. The pilot scored 47/60 with no operational failures, so this model entered the frozen benchmark. \\
Qwen3.8-27B-FP8 & A 32,768-token screen solved easy but exhausted context on medium and hard. A larger-context medium rerun was correct. This screen is separate from the frozen benchmark. \\
Qwen3.5-122B-A10B FP8 & The model was prepared, but its diagnostic job was cancelled before it started. We draw no accuracy conclusion. \\
Mistral Small 4 119B and Mistral Medium 3.5 128B & These models were prepared or configured but had no completed scored screen. We omitted them for scope and compute control. \\
\bottomrule
\end{tabularx}
\end{table}

Other open-weight families, including GLM and Kimi, were discussed during candidate search but were not run as scored experiments. The resulting thesis is a controlled single-model representation study. The candidate screens informed model selection only. They did not alter the frozen puzzle samples or the main request set.

\subsubsection{Representations}
The main benchmark uses nine ordered alphabets.  Each item in an alphabet represents the abstract value in the same position:
\begin{center}
\small
\begin{tabularx}{\linewidth}{@{}lX@{}}
\toprule
Family & Ordered labels for $V_1$ through $V_9$ \\
\midrule
Arabic digits & 1, 2, 3, 4, 5, 6, 7, 8, 9 \\
Devanagari numerals & {\devanagarifont १, २, ३, ४, ५, ६, ७, ८, ९} \\
Bengali numerals & {\banglafont ১, ২, ৩, ৪, ৫, ৬, ৭, ৮, ৯} \\
Uppercase Latin & A, B, C, D, E, F, G, H, I \\
Lowercase Latin & a, b, c, d, e, f, g, h, i \\
Greek letters & $\alpha$, $\beta$, $\gamma$, $\delta$, $\varepsilon$, $\zeta$, $\eta$, $\theta$, $\iota$ \\
Abstract symbols & triangle, square, circle, star, multiplication sign, dagger, double dagger, section sign, paragraph sign \\
Emoji & red, orange, yellow, green, blue, purple, brown, black, and white circles \\
Nonce labels & KAV, MIP, ZOT, RUL, BEK, DAX, PEV, NUG, WIF \\
\bottomrule
\end{tabularx}
\end{center}
All symbols are stored in Unicode NFC form.  The encoder writes UTF-8 text, and a round-trip test requires decoding an encoded grid to recover the original abstract grid exactly.  The dot is the default empty-cell marker and cannot be a value symbol.

\subsubsection{Standard prompt and response}
The standard prompt names the valid alphabet, shows the transformed puzzle, and states four requirements: every row, column, and $3\times3$ box must contain each value exactly once, and the given cells must not change.  It also states that visible symbols are labels and that one mapping must be used consistently.  The requested answer is exactly nine lines of nine single-space-separated output symbols, with no explanation.  Apart from the declared alphabet and transformed cell values, the puzzle, wording, output format, model settings, and evaluator remain fixed.

The exact parser does not search prose for a plausible grid.  It requires nine rows, 81 cells, one space between adjacent symbols, and membership in the requested alphabet.  A candidate is correct only if it is parseable, preserves every clue, equals the certified solution, and satisfies all nine rows, nine columns, and nine boxes.  The evaluator can attach several error labels to one answer: changed clue, row, column or box violation, invalid symbol, wrong row or cell count, format error, Unicode error, refusal, missing final answer, or truncation.  If one consistent permutation changes the returned grid into the certified solution, it is a global binding error.  Other inconsistent cell mappings are local mapping errors.

An API error or timeout produces \texttt{NOT\_EVALUATED}.  It is reported as an operational failure and is not counted as an incorrect Sudoku answer.  An answer that finishes normally but is invalid, malformed, missing, refused, or truncated is an evaluated model failure.  Backend failures follow the fixed retry policy of at most five attempts with exponential delays.  The raw status, token counts, finish reason, timestamps, and generation latency are saved with every result.

\subsection{Experiment 1: dataset and evaluator validation}
\paragraph{Question.}
Can later failures be attributed to the model rather than to an ambiguous puzzle, a wrong reference solution, or a Unicode and parser defect?

\paragraph{Method.}
A seeded exact solver first generates a completed grid.  Clues are then considered for removal in seeded random order.  A removal is kept only when an exact solution counter, which stops after two solutions, finds exactly one solution.  The generator targets 40--46 clues for easy candidates, 24--27 for medium candidates, and 22--26 for hard candidates.  Candidates whose measured difficulty does not match the requested tier are rejected, as are duplicate puzzle hashes.

Difficulty is assigned by a deterministic technique solver.  A puzzle is easy if naked and hidden singles solve it.  It is medium if solving also needs an implemented intermediate technique: naked pair, hidden pair, pointing pair, or box-line reduction.  If these techniques stop before completion and exact search is required, the puzzle is hard.  The certificate records deduction steps, techniques used, whether backtracking was required, maximum search depth, and search nodes.

The final dataset contains 300 unique puzzles: 100 easy, 100 medium, and 100 hard.  Each record stores its ID, compact puzzle and solution, clue mask and count, generation seed and version, SHA-256 puzzle hash, difficulty certificate, and solver certificate.  A separate audit rechecks record structure, hashes, uniqueness, exact solution count, stored solution, difficulty, and certificate.  File hashes are also stored in a manifest.  This experiment therefore establishes both valid test items and a deterministic scoring pipeline before any model comparison.

\subsection{Experiment 2: model screen, qualification, and pilot}
\paragraph{Question.}
Can the selected model and compute setup complete exact Sudoku grids across difficulty levels and transformed alphabets, and are the selected puzzles informative rather than all too easy or all too hard?

\paragraph{Method.}
A timing screen uses one reproducibly selected puzzle from each difficulty.  Its purpose is to measure natural reasoning time and to set scheduler wall time without limiting reasoning to a small token budget.  Qualification then uses five separate easy puzzles.  They are assigned once each to Arabic digits, Greek letters, emoji, nonce labels, and Devanagari numerals.  The model must solve all five with no operational failure before the protocol can be frozen.

The pilot uses 15 different puzzles, with five from each difficulty tier.  Every puzzle is shown under Arabic digits, uppercase Latin, Greek, and emoji, giving $15\times4=60$ requests.  The same puzzle set in every representation makes the representation comparison balanced.  Pilot accuracy is inspected by difficulty and representation, together with latency, output length, truncation, and operational failures.  Pilot puzzles are excluded from the main sample.  Thus, any choice informed by pilot performance cannot change the confirmatory outcomes for those same items.

\subsection{Experiment 3: representation registry and protocol freeze}
\paragraph{Question.}
Can a representation effect be related to measurable symbol cost, and can the confirmatory design be protected from changes made after seeing results?

\paragraph{Method.}
For every main-alphabet symbol, the registry records its Unicode code points, UTF-8 byte length, code-point count, grapheme count, token count and token IDs in isolation, token count after a leading space, and the token cost of the full alphabet row.  For every puzzle and alphabet it also records clue count, total tokens used by visible clues, mean clue-token cost, and total prompt tokens.  Exact token IDs are available because the selected model has a local tokenizer.

After qualification and the pilot, the code freezes the dataset hash, master seed, sample IDs, model revision, inference and retry settings, experiment shard counts, and a digest of all numbered scripts and library source files.  It also freezes the exact request-ID set for each experiment.  Later scripts refuse to run if the configuration, sample plan, or source digest differs from the frozen protocol.

Sample selection is deterministic and stratified.  The 15-puzzle pilot is disjoint from the main sample.  The main sample contains 60 puzzles, with 20 per difficulty.  A 15-puzzle mechanism sample, with five per difficulty, is selected from the main sample.  A nine-puzzle ablation sample, with three per difficulty, is selected from the mechanism sample.  This nesting allows later experiments to reuse the exact main-benchmark answers instead of generating a different baseline.

\subsection{Experiment 4: main representation benchmark}
\paragraph{Question.}
Does changing only the visible value labels change complete-grid accuracy?

\paragraph{Method.}
Every one of the 60 main puzzles is presented under all nine alphabets.  This gives $60\times9=540$ requests for the selected model.  All logical content and inference settings are held fixed.  The primary summary for each family is the number correct divided by the number evaluated, with operational failures shown separately and Wilson 95\% confidence intervals for proportions.  Results are also divided by difficulty, and diagnostic failure labels are counted.

For reliable long runs, request IDs are SHA-256 hashes of canonical request data.  A second hash assigns the 540 requests to six deterministic shards.  Results are appended to JSONL after each request.  A resumed run skips existing request IDs, and a shard is complete only when all eligible requests are present and a completion marker has been written.  These controls prevent a scheduler timeout from losing finished work or causing duplicate inference.

\subsection{Experiment 5: paired Arabic retention analysis}
\paragraph{Question.}
Among puzzles the model can solve with ordinary digits, how often does it retain that ability after the values are renamed?

\paragraph{Method.}
This experiment makes no new model calls.  For each non-Arabic family $s$, let $B$ be the main-sample puzzles answered correctly in Arabic.  The retention and representation-failure rates are
\begin{equation}
R_s=\frac{|\{p\in B:p\text{ is correct under }s\}|}{|B|},
\qquad F_s=1-R_s.
\end{equation}
They are reported overall and by difficulty.  This conditional analysis directly answers the thesis question, because it asks whether a demonstrated Arabic solution survives the transformation.  All paired Arabic-only and transformed-only successes are also compared with a two-sided exact McNemar test.  Cell, constraint, binding, termination, and operational failures remain separate.

\subsection{Experiment 6: input/output cross and basic symbol controls}
\paragraph{Question.}
If transformed performance is lower, is the main problem reading the transformed puzzle, solving while maintaining the mapping, or producing a long transformed answer?

\paragraph{Method.}
The 15 mechanism puzzles are tested in four conditions. They use Arabic input and output (A), Greek input and output (B), Greek input with Arabic output (C), and Arabic input with Greek output (D). Cross-alphabet prompts give the complete symbol mapping. Conditions A and B reuse matching main-benchmark results. Conditions C and D require new calls. The comparisons help distinguish a transformed-input cost from a transformed-output cost, although they cannot isolate every stage of a long Sudoku solution.

Five no-Sudoku controls are added for every mechanism puzzle. They require copying a Greek row, translating a short sequence, retrieving one coordinate, counting one symbol, or converting a supplied complete grid. They use exact-text scoring. These controls test whether elementary symbol operations fail without Sudoku search. The experiment contains 135 recorded conditions. Thirty Sudoku baselines are reused, and 105 requests require new inference.

\subsection{Experiment 7: token-length mechanism}
\paragraph{Question.}
Does accuracy decrease because some visible labels require more model tokens, rather than because their abstract meaning is different?

\paragraph{Method.}
Using the selected model's tokenizer and the master seed, the code searches neutral uppercase nonce strings.  It constructs three nine-label alphabets whose labels require exactly one, two, or three tokens both in isolation and after a leading space.  The 15 mechanism puzzles are solved under all three alphabets, giving 45 requests.  If any token-length group cannot supply nine verified distinct labels, the experiment stops instead of silently changing its design.

The analysis relates correctness to total prompt tokens, mean clue tokens, mean UTF-8 bytes, and mean code points using logistic regression with standardized predictors.  Token length is supported as a mechanism only if accuracy changes in the controlled alphabets.  A remaining family effect would show that token count is incomplete, not that tokenization is irrelevant.

\subsection{Experiment 8: arbitrary binding and semantic conflict}
\paragraph{Question.}
Is performance sensitive to an arbitrary mapping itself, and do familiar labels interfere when they are assigned meanings that conflict with their usual meanings?

\paragraph{Method.}
The 15 mechanism puzzles use the standard uppercase mapping and five deterministic random permutations of the same letters.  Variation across these five mappings measures sensitivity to the particular binding rather than to the alphabet alone.  The same puzzles also use ordinary digits, a fixed permutation of digits, ordinary English number words, a fixed conflicting permutation of number words, and neutral nonce labels.  A larger penalty for conflicting familiar labels than for neutral labels supports semantic interference.

There are 11 conditions per puzzle, or 165 recorded conditions.  Standard uppercase, ordinary digits, and nonce labels reuse matching main-benchmark results.  The remaining eight conditions per puzzle require 120 new calls.  Seed-fixed permutations make the comparison reproducible.

\subsection{Experiment 9: prompt and format ablations}
\paragraph{Question.}
Could the measured representation effect be reduced or explained by more explicit instructions, a declared mapping, a different output format, a different empty marker, or case and spacing conventions?

\paragraph{Method.}
This is an exploratory one-factor-at-a-time study on the nine-puzzle ablation sample.  Greek puzzles are tested with four rule styles: minimal, explicit constraints, constraints plus alphabet, and the fully explicit standard prompt.  Three mapping styles provide only the alphabet, a mapping to digits, or a mapping to abstract values.  Four output formats require spaced rows, compact rows, one 81-symbol string, or a JSON array.  Four empty markers are compared: dot, zero, underscore, and the word \texttt{EMPTY}.  Finally, uppercase and lowercase forms of Latin and nonce labels are compared.

These 19 variants give $9\times19=171$ requests.  Each output format has its own exact normalizer before the common Sudoku evaluator.  Ablation results do not retroactively change the frozen main protocol.

\subsection{Experiment 10: self-revision and checker-guided revision}
\paragraph{Question.}
Can the model detect and repair an incorrect grid, and is exact checker feedback more useful than an unaided request to check again?

\paragraph{Method.}
The nine ablation puzzles are used in Arabic, Greek, and emoji form.  For each puzzle-representation pair, four independent branches start from the same saved main-benchmark answer: no revision, one self-revision, two self-revisions, and one checker-guided revision.  Reusing the same initial answer removes variation caused by generating four different baselines.  Each branch has its own conversation after that initial answer.

The self-revision prompt asks the model to check cell count, clues, rows, columns, and boxes.  In the checker-guided branch, the deterministic evaluator returns error descriptions, such as a changed clue or a duplicate and missing symbol, but never supplies the correct cell value.  A correct initial answer receives no unnecessary checker-guided call.  Every intermediate and final grid is saved and scored.  The analysis reports initial and final accuracy, wrong answers fixed, correct answers made wrong, model calls, token use, and total generation time.  The 108 branch records require between 81 and 108 new calls, depending on how many checker-guided baselines are already correct.

\subsection{Analysis and reporting}
The unit of correctness is the complete final grid. We do not use partial credit or visible reasoning quality. Accuracy summaries exclude only \texttt{NOT\_EVALUATED} operational cases from their denominator and report those cases separately. Counts accompany percentages so that small groups remain visible. Wilson 95\% intervals summarize accuracy proportions. Exact McNemar tests compare paired main outcomes, and the token mechanism uses the regression described above. We also report format compliance, clue preservation, constraint validity, mapping errors, truncation, latency, token counts, and operational status. This keeps invalid grids and transport failures separate from successful solutions.

\section{Results}
The main benchmark addresses RQ1. The cross-mapping, token-length, and binding experiments address RQ2. The prompt ablations and self-revision experiment address RQ3. Qualification and pilot results establish that the chosen model and prompts could produce exact grids before the main protocol was frozen.
\subsection{Qualification and timing}
GPT-OSS-120B solved all five qualification requests and produced no operational failures. Mean qualification latency was \SI{44.6}{\second} per request. In the one-puzzle timing screen, the easy, medium, and hard puzzles were all solved correctly in \SI{36.1}{\second}, \SI{254.6}{\second}, and \SI{227.4}{\second}, respectively. These measurements motivated wall-time-based execution rather than a small reasoning-token budget.

\subsection{Balanced pilot}
The model solved \PilotCorrect{} pilot requests (\PilotAccuracy{}). All 60 requests completed at the inference layer and none had an operational error. Mean generation latency was \PilotMeanLatency{} seconds, the median was \PilotMedianLatency{} seconds, and mean completion length was \PilotMeanTokens{} tokens. Two responses reached the generation ceiling and did not contain a complete final grid.

\begin{table}[H]
\centering
\caption{Pilot accuracy by difficulty. Intervals are Wilson 95\% confidence intervals.}
\label{tab:pilot-difficulty}
\begin{tabular}{lrrr}
\toprule
Difficulty & Correct & Accuracy (\%) & 95\% CI (\%) \\
\midrule
Easy & 20/20 & 100.0 & [83.9, 100.0] \\
Hard & 14/20 & 70.0 & [48.1, 85.5] \\
Medium & 13/20 & 65.0 & [43.3, 81.9] \\
\bottomrule
%
\end{tabular}
\end{table}

\begin{figure}[H]
\centering
\begin{tikzpicture}
\begin{axis}[
  xbar, xmin=0, xmax=105, width=0.82\linewidth, height=5.6cm,
  xlabel={Accuracy (\%)}, symbolic y coords={medium,hard,easy},
  ytick=data, bar width=16pt, nodes near coords,
  nodes near coords style={font=\small,/pgf/number format/fixed,/pgf/number format/precision=1},
  axis line style={draw=black!55}, tick style={draw=black!55},
  ymajorgrids=true, grid style={black!10},
]
\addplot+[fill=accent,draw=accent] table[x=accuracy,y=label] {generated/pilot_difficulty.dat};
\end{axis}
\end{tikzpicture}
\caption{Pilot accuracy by difficulty ($n=20$ per tier).}
\label{fig:pilot-difficulty}
\end{figure}

\begin{table}[H]
\centering
\caption{Pilot accuracy by representation.}
\label{tab:pilot-representation}
\begin{tabular}{lrrr}
\toprule
Representation & Correct & Accuracy (\%) & 95\% CI (\%) \\
\midrule
Arabic digits & 14/15 & 93.3 & [70.2, 98.8] \\
Emoji & 13/15 & 86.7 & [62.1, 96.3] \\
Greek letters & 10/15 & 66.7 & [41.7, 84.8] \\
Uppercase Latin & 10/15 & 66.7 & [41.7, 84.8] \\
\bottomrule
%
\end{tabular}
\end{table}

\begin{figure}[H]
\centering
\begin{tikzpicture}
\begin{axis}[
  xbar, xmin=0, xmax=105, width=0.88\linewidth, height=6.2cm,
  xlabel={Accuracy (\%)}, symbolic y coords={uppercase-latin,greek-letters,emoji,arabic-digits},
  ytick=data, bar width=13pt, nodes near coords,
  nodes near coords style={font=\small,/pgf/number format/fixed,/pgf/number format/precision=1},
  axis line style={draw=black!55}, tick style={draw=black!55},
  xmajorgrids=true, grid style={black!10},
]
\addplot+[fill=accentlight,draw=accent] table[x=accuracy,y=label] {generated/pilot_representation.dat};
\end{axis}
\end{tikzpicture}
\caption{Pilot accuracy by representation ($n=15$ per representation).}
\label{fig:pilot-representation}
\end{figure}

Table~\ref{tab:pilot-difficulty} and Figure~\ref{fig:pilot-difficulty} show that all 20 easy pilot requests were correct. Medium and hard accuracy was 13/20 and 14/20. Table~\ref{tab:pilot-representation} and Figure~\ref{fig:pilot-representation} show 14/15 for Arabic digits, 13/15 for emoji, and 10/15 each for Greek and uppercase Latin. Every representation used the same 15 puzzles, so difficulty composition did not cause these differences. The medium-hard ordering was not monotonic in this small pilot.

\subsection{Main representation benchmark}
All six deterministic shards completed, giving 540/540 recorded requests and valid completion markers. GPT-OSS-120B solved \MainCorrect{} requests (\MainAccuracy{}) with \MainOperationalFailures{} operational failures. Mean and median generation latency were \MainMeanLatency{} and \MainMedianLatency{} seconds, and mean completion length was \MainMeanTokens{} tokens. Twenty-four responses reached the generation ceiling. Exact row formatting was preserved in \MainFormatCompliant{}/540 responses. Of the \MainParseable{} responses that could be decoded as grids over the requested alphabet, \MainCluesPreserved{} preserved every given cell.

\begin{table}[H]
\centering
\caption{Main-benchmark accuracy by difficulty. Each tier contains 20 puzzles under nine representations.}
\label{tab:main-difficulty}
\begin{tabular}{lrrr}
\toprule
Difficulty & Correct & Accuracy (\%) & 95\% CI (\%) \\
\midrule
Easy & 171/180 & 95.0 & [90.8, 97.3] \\
Hard & 80/180 & 44.4 & [37.4, 51.7] \\
Medium & 121/180 & 67.2 & [60.1, 73.7] \\
\bottomrule
%
\end{tabular}
\end{table}

\begin{table}[H]
\centering
\caption{Main-benchmark accuracy by representation. Each representation uses the same 60 puzzles.}
\label{tab:main-representation}
\begin{tabular}{lrrr}
\toprule
Representation & Correct & Accuracy (\%) & 95\% CI (\%) \\
\midrule
Abstract symbols & 38/60 & 63.3 & [50.7, 74.4] \\
Arabic digits & 45/60 & 75.0 & [62.8, 84.2] \\
Bengali numerals & 43/60 & 71.7 & [59.2, 81.5] \\
Devanagari numerals & 39/60 & 65.0 & [52.4, 75.8] \\
Emoji & 38/60 & 63.3 & [50.7, 74.4] \\
Greek letters & 37/60 & 61.7 & [49.0, 72.9] \\
Lowercase Latin & 44/60 & 73.3 & [61.0, 82.9] \\
Nonce labels & 43/60 & 71.7 & [59.2, 81.5] \\
Uppercase Latin & 45/60 & 75.0 & [62.8, 84.2] \\
\bottomrule
%
\end{tabular}
\end{table}

\begin{figure}[H]
\centering
\begin{tikzpicture}
\begin{axis}[
  xbar, xmin=0, xmax=100, width=0.84\linewidth, height=8.4cm,
  xlabel={Accuracy (\%)},
  symbolic y coords={greek-letters,emoji,abstract-symbols,devanagari-numerals,bengali-numerals,nonce-labels,lowercase-latin,uppercase-latin,arabic-digits},
  ytick=data, bar width=10pt, nodes near coords,
  nodes near coords style={font=\scriptsize,/pgf/number format/fixed,/pgf/number format/precision=1},
  axis line style={draw=black!55}, tick style={draw=black!55},
  xmajorgrids=true, grid style={black!10},
]
\addplot+[fill=accent,draw=accent] table[x=accuracy,y=label] {generated/main_representation.dat};
\end{axis}
\end{tikzpicture}
\caption{Main-benchmark accuracy by representation ($n=60$ each).}
\label{fig:main-representation}
\end{figure}

RQ1 concerns representation and difficulty on the fixed main sample. Table~\ref{tab:main-difficulty} shows 171/180 correct on easy, 121/180 on medium, and 80/180 on hard requests. Table~\ref{tab:main-representation} and Figure~\ref{fig:main-representation} show a range from 37/60 for Greek letters to 45/60 for Arabic digits and uppercase Latin. Equal total accuracy for two alphabets does not imply that they succeeded on the same puzzles.

\subsection{Paired Arabic retention}
Arabic digits were correct for 45 of the 60 main puzzles. Table~\ref{tab:retention} conditions each transformed representation on those 45 demonstrated Arabic successes. Retention ranged from 31/45 (68.9\%) for abstract symbols to 37/45 (82.2\%) for lowercase Latin. The discordant columns report Arabic-only/transformed-only successes across all 60 paired puzzles.

\begin{table}[H]
\centering
\small
\caption{Retention among the 45 puzzles solved correctly with Arabic digits. McNemar $p$ values are exact and two-sided.}
\label{tab:retention}
\begin{tabular}{lrrrr}
\toprule
Representation & Retained & Retention (\%) & Discordant & $p$ \\
\midrule
Abstract symbols & 31/45 & 68.9 & 14/7 & 0.189 \\
Bengali numerals & 36/45 & 80.0 & 9/7 & 0.804 \\
Devanagari numerals & 36/45 & 80.0 & 9/3 & 0.146 \\
Emoji & 32/45 & 71.1 & 13/6 & 0.167 \\
Greek letters & 36/45 & 80.0 & 9/1 & 0.021 \\
Lowercase Latin & 37/45 & 82.2 & 8/7 & 1.000 \\
Nonce labels & 36/45 & 80.0 & 9/7 & 0.804 \\
Uppercase Latin & 36/45 & 80.0 & 9/9 & 1.000 \\
\bottomrule
%
\end{tabular}
\end{table}

\subsection{Input/output cross and symbol controls}
The input/output-cross experiment completed all 135 recorded conditions with \InputOutputOperationalFailures{} operational failures. Of these, 30 baseline Sudoku results were reused from the main benchmark and 105 were new inference calls. Table~\ref{tab:input-output} reports every condition.

\begin{table}[H]
\centering
\small
\caption{Input/output-cross and no-Sudoku control accuracy ($n=15$ per condition).}
\label{tab:input-output}
\begin{tabular}{lrrr}
\toprule
Condition & Correct & Accuracy (\%) & 95\% CI (\%) \\
\midrule
A: Arabic input, Arabic output & 12/15 & 80.0 & [54.8, 93.0] \\
B: Greek input, Greek output & 12/15 & 80.0 & [54.8, 93.0] \\
C: Greek input, Arabic output & 10/15 & 66.7 & [41.7, 84.8] \\
D: Arabic input, Greek output & 11/15 & 73.3 & [48.0, 89.1] \\
Coordinate retrieval & 15/15 & 100.0 & [79.6, 100.0] \\
Greek row copy & 15/15 & 100.0 & [79.6, 100.0] \\
Grid conversion & 15/15 & 100.0 & [79.6, 100.0] \\
Mapping translation & 15/15 & 100.0 & [79.6, 100.0] \\
Occurrence count & 15/15 & 100.0 & [79.6, 100.0] \\
\bottomrule
%
\end{tabular}
\end{table}

All five basic symbol controls were correct for all 15 puzzles: copying, short mapping translation, coordinate retrieval, occurrence counting, and completed-grid conversion each scored 15/15. The four Sudoku conditions were lower: Arabic input to Arabic output and Greek input to Greek output each scored 12/15, Greek input to Arabic output scored 10/15, and Arabic input to Greek output scored 11/15. By difficulty across the four Sudoku conditions, accuracy was 19/20 on easy, 18/20 on medium, and 8/20 on hard puzzles. The two truncated responses occurred on hard Sudoku requests.

\subsection{Token-length mechanism}
The token-length experiment completed all 45 requests with \TokenLengthOperationalFailures{} operational failures and \TokenLengthTruncations{} truncated responses. Accuracy was not monotonic in the controlled label length.

\begin{table}[H]
\centering
\caption{Accuracy under neutral alphabets requiring one, two, or three tokens per label ($n=15$ each).}
\label{tab:token-length}
\begin{tabular}{lrrr}
\toprule
Condition & Correct & Accuracy (\%) & 95\% CI (\%) \\
\midrule
One-token labels & 9/15 & 60.0 & [35.7, 80.2] \\
Two-token labels & 13/15 & 86.7 & [62.1, 96.3] \\
Three-token labels & 10/15 & 66.7 & [41.7, 84.8] \\
\bottomrule
%
\end{tabular}
\end{table}

The one-token labels scored 9/15 (60.0\%), the two-token labels scored 13/15 (86.7\%), and the three-token labels scored 10/15 (66.7\%). By difficulty across the three alphabets, accuracy was 15/15 on easy, 13/15 on medium, and 4/15 on hard puzzles. The five truncations comprised four hard requests and one medium request. The preregistered multivariable logistic fit could not be estimated because its information matrix was singular: prompt length, mean clue-token cost, byte length, and code-point length co-varied deterministically in this 45-request construction. Accordingly, the descriptive result rejects a simple monotonic token-count account but does not identify an independent coefficient for token length.

\subsection{Arbitrary binding and semantic conflict}
The binding experiment completed all 165 records with \BindingOperationalFailures{} request-level operational failures. It solved \BindingCorrect{} conditions (\BindingAccuracy{}) and produced \BindingTruncations{} truncated responses. Table~\ref{tab:binding} reports all mappings.

\begin{table}[H]
\centering
\small
\caption{Binding accuracy and truncation by condition ($n=15$ each).}
\label{tab:binding}
\begin{tabular}{lrrr}
\toprule
Condition & Correct & Accuracy (\%) & Truncated \\
\midrule
Standard uppercase & 11/15 & 73.3 & 1 \\
Uppercase permutation 1 & 11/15 & 73.3 & 1 \\
Uppercase permutation 2 & 11/15 & 73.3 & 1 \\
Uppercase permutation 3 & 13/15 & 86.7 & 0 \\
Uppercase permutation 4 & 12/15 & 80.0 & 0 \\
Uppercase permutation 5 & 10/15 & 66.7 & 0 \\
Ordinary digits & 12/15 & 80.0 & 1 \\
Permuted digits & 10/15 & 66.7 & 1 \\
Ordinary number words & 12/15 & 80.0 & 0 \\
Conflicting number words & 10/15 & 66.7 & 1 \\
Neutral nonce labels & 12/15 & 80.0 & 1 \\
\bottomrule
%
\end{tabular}
\end{table}

All 55 easy conditions were correct. Medium and hard accuracy was 50/55 and 19/55. The six mappings over the identical uppercase symbol set ranged from 10/15 to 13/15, with the standard mapping at 11/15. Ordinary digits, ordinary number words, and neutral nonce labels each scored 12/15. Permuted digits and conflicting number words each scored 10/15.

\begin{table}[H]
\centering
\small
\caption{Paired binding outcomes. Discordance is first-condition-only/second-condition-only success.}
\label{tab:binding-pairs}
\begin{tabular}{lrrr}
\toprule
Comparison & First & Second & Discordant \\
\midrule
Digits & 12/15 & 10/15 & 2/0 \\
Number words & 12/15 & 10/15 & 3/1 \\
Nonce/conflicting words & 12/15 & 10/15 & 2/0 \\
\bottomrule
%
\end{tabular}
\end{table}

For ordinary versus permuted digits, two puzzles were solved only under the ordinary mapping and none only under the permutation. Ordinary versus conflicting number words had three ordinary-only and one conflict-only success. Neutral nonce labels versus conflicting words had two neutral-only and no conflict-only success. These directional differences are small relative to the 15-puzzle sample. They provide descriptive evidence of semantic interference, not a conclusive effect. Equal truncation counts in the digit and nonce/conflicting comparisons show that differential truncation did not produce their two-answer gaps.

\subsection{Prompt and output ablations}
The exploratory ablation experiment recorded all 171 planned conditions and its completion marker before Slurm terminated backend teardown at the 15-hour wall-time. The request records contain \AblationOperationalFailures{} operational failures. Overall accuracy was \AblationCorrect{} (\AblationAccuracy{}), with \AblationTruncations{} truncated responses. Table~\ref{tab:ablations} reports every variant.

\begin{table}[H]
\centering
\small
\caption{Exploratory prompt and output ablations ($n=9$ per variant).}
\label{tab:ablations}
\begin{tabular}{llrr}
\toprule
Factor & Variant & Correct & Accuracy (\%) \\
\midrule
Rules & Minimal & 5/9 & 55.6 \\
Rules & Explicit constraints & 6/9 & 66.7 \\
Rules & Constraints and alphabet & 5/9 & 55.6 \\
Rules & Fully explicit & 5/9 & 55.6 \\
\addlinespace
Mapping & Alphabet only & 6/9 & 66.7 \\
Mapping & To digits & 8/9 & 88.9 \\
Mapping & To abstract values & 5/9 & 55.6 \\
\addlinespace
Output & Spaced rows & 5/9 & 55.6 \\
Output & Compact rows & 3/9 & 33.3 \\
Output & 81-symbol string & 3/9 & 33.3 \\
Output & JSON & 5/9 & 55.6 \\
\addlinespace
Empty marker & Dot & 7/9 & 77.8 \\
Empty marker & Zero & 6/9 & 66.7 \\
Empty marker & Underscore & 4/9 & 44.4 \\
Empty marker & \texttt{EMPTY} & 6/9 & 66.7 \\
\addlinespace
Latin case & Uppercase & 5/9 & 55.6 \\
Latin case & Lowercase & 7/9 & 77.8 \\
\addlinespace
Nonce case & Uppercase & 8/9 & 88.9 \\
Nonce case & Lowercase & 4/9 & 44.4 \\
\bottomrule
%
\end{tabular}
\end{table}

Rule wording ranged only from 5/9 to 6/9 and showed no advantage for the fully explicit prompt. Declaring a digit mapping scored 8/9, compared with 6/9 for alphabet-only and 5/9 for abstract-value mapping. Neither compact rows nor an 81-symbol string improved on spaced rows or JSON. Empty-marker results ranged from 4/9 for underscore to 7/9 for dot. Case effects reversed across families: lowercase Latin exceeded uppercase Latin (7/9 versus 5/9), whereas uppercase nonce labels exceeded lowercase nonce labels (8/9 versus 4/9).

Four named variants produce identical prompts for each puzzle under the one-factor construction. They are fully explicit rules, alphabet-only mapping, spaced output, and dot empty marker. Their independent sampled runs scored 5/9, 6/9, 5/9, and 7/9. This two-answer spread among exact prompt replicates shows sampling variability at this scale. The larger ablation contrasts are useful diagnostics, but none warrants changing the frozen main protocol on nine puzzles alone.

\subsection{Self-revision and checker-guided revision}
The answer-revision experiment addresses RQ3 on 27 puzzle-representation pairs. It saved all 108 branch records with no request-level operational failures. Table~\ref{tab:revision} compares the final answer in each branch. A correct initial answer required no checker-guided call.

\begin{table}[H]
\centering
\caption{Complete-grid accuracy after answer revision ($n=27$ per branch).}
\label{tab:revision}
\begin{tabular}{lrr}
\toprule
Branch & Correct & Accuracy (\%) \\
\midrule
One pass & 18/27 & 66.7 \\
One self-revision & 19/27 & 70.4 \\
Two self-revisions & 23/27 & 85.2 \\
Checker-guided revision & 22/27 & 81.5 \\
\bottomrule
\end{tabular}
\end{table}

Two self-revisions produced five more correct final grids than one pass. Checker-guided revision produced four more. These are outcomes for four branches from the same saved initial answers, but each branch makes different model calls. The 27-item sample supports an exploratory improvement on these puzzles. It does not establish that two revisions are generally best or that model self-checking is equivalent to deterministic verification.

\section{Discussion}
The answer to RQ1 is that GPT-OSS-120B's complete-grid accuracy changed across representations of the same Sudoku puzzles. The gap between Greek letters and Arabic digits was 13.3 percentage points across 60 paired puzzles. Difficulty had a larger association with accuracy. The model solved 95.0\% of easy requests, 67.2\% of medium requests, and 44.4\% of hard requests.

Marginal accuracy hides item-level changes. Arabic digits and uppercase Latin letters both scored 45/60. Yet uppercase Latin retained only 36 of the 45 Arabic successes, while nine other puzzles changed from Arabic failure to uppercase-Latin success. Thus, equal totals did not mean that the two representations supported the same solutions. Bengali numerals also matched ASCII nonce labels at 71.7\%. These outcomes do not follow a simple ASCII-versus-Unicode ordering.

RQ2 remains only partly answered. The model completed all 75 basic symbol controls, including copying and conversion of a supplied complete grid. It therefore handled those isolated tasks under the tested prompts. Cross-mapping Sudoku accuracy was lower than both same-alphabet conditions on the 15-puzzle mechanism sample. The sample does not separate input decoding, mapping maintenance, and output generation into independent effects.

Token length was not monotonically associated with accuracy in the controlled nonce-label test. Two-token labels performed best, and the planned regression could not identify separate coefficients because the predictors co-varied. Familiar-label conflicts scored two fewer correct answers than their ordinary counterparts in the 15-puzzle binding sample. Several mappings of the same uppercase tokens also gave different outcomes. Together, these tests suggest that binding and familiar meanings may matter, but they do not establish their individual contributions.

RQ3 also produced mixed evidence. No prompt or output variant consistently improved accuracy in the nine-puzzle ablation. Exact prompt replicates varied by two correct answers, which limits interpretation of small differences. On 27 puzzle-representation pairs, two self-revisions yielded 23 correct final grids, compared with 18 after one pass. Checker-guided revision yielded 22. These exploratory gains warrant testing on a larger sample before changing the frozen benchmark procedure.

Generation cost was substantial. Main-benchmark requests averaged \MainMeanLatency{} seconds and \MainMeanTokens{} completion tokens. Twenty-four responses reached the generation ceiling. Every main request finished at the inference layer, so these truncated answers were model-output failures rather than operational errors. The study kept the ceiling fixed across main representations to preserve comparability.

\section{Limitations}
The main benchmark evaluates one model checkpoint and one prompt protocol. Its 60 puzzles were generated and certified by the study's solver, so the findings need not transfer to human-authored Sudoku collections or other reasoning tasks. Each representation appears once per puzzle. Sampling variability is visible in the exact-prompt ablation replicates, but the main benchmark does not estimate repeated-generation variance for every condition. The smaller mechanism and revision samples support descriptive comparisons, not precise estimates of causal effects. We score only the final grid and make no claims about the correctness of hidden reasoning.

\section{Conclusion}
We built an exact, paired test of Sudoku solving under nine value alphabets. GPT-OSS-120B solved \MainCorrect{} frozen benchmark requests with no operational failures. Its accuracy ranged from 61.7\% to 75.0\% across representations of the same 60 puzzles. The paired analysis showed that even equal total accuracy can conceal different item-level successes. Basic symbol controls were all correct, while token-length, binding, prompt, and revision tests gave mixed or sample-limited results. The measured Sudoku performance was sensitive to value renaming, and the completed experiments do not identify a single sufficient explanation.

\begingroup
\small
\setlength{\bibsep}{2pt}
\bibliographystyle{plainnat}
\bibliography{references}
\endgroup

\end{document}
